\section{Compact charge constraints}
\label{sec:compact-gauging}\label{sec:local-sectors}

\subsection{Wrapped-brane charges and deformation coordinates}
\label{subsec:m2-charges}

Let $\pi:\wideX\to X$ be a projective crepant resolution of a
compact Calabi--Yau threefold with isolated singularities. In an
integral divisor basis, the three-form reduction and wrapped-M2
coupling are
\begin{equation}
 C_3=\sum_I A^I\wedge\omega_I+\cdots,\qquad
 q_I(C)=\int_C\omega_I=D_I C.
 \label{eq:c3-reduction}
\end{equation}
For a resolved node the exceptional curve supports the conifold
hypermultiplet \cite{Greene:1995hu,Katz:1996,Mohaupt:2004de}.
In one complex structure its scalars $(h,\widetilde h)$ have
opposite charges, and the invariant deformation coordinate is
\begin{equation}
 u=h\widetilde h,\qquad xy-zw=u.
 \label{eq:conifold-map}
\end{equation}
For $u=|u|e^{i\theta}$, taking
$(h,\widetilde h)=\sqrt{|u|}(e^{i\theta},1)$ also sets the
real moment map to zero.

The anticanonical cone over $S_7=\operatorname{Bl}_2\mathbb P^2$
engineers the rank-one $E_2$ theory, with flavor algebra
$\mathfrak{su}(2)\oplus\mathfrak u(1)$
\cite{Seiberg:1996,Morrison:1996pp}. Its reduced Higgs ring is
\begin{equation}
 \CC[J_+,J_0,J_-]/(J_+J_--J_0^2),\qquad \beta=J_0.
 \label{eq:e2-reduced}
\end{equation}
The operators $J_\pm$ are the instanton currents in its
$SU(2)$ gauge-theory deformation with one flavor
\cite{Cremonesi:2015lsa}. The geometric deformation coordinate
$\beta$ is the Cartan invariant, not the whole two-complex-dimensional
Higgs cone. Normalize its effective Cartan action to give $J_+$
charge one. On the double cover,
\begin{equation}
 (J_+,J_0,J_-)=(a^2,ab,b^2)/2,\quad(a,b)\sim(-a,-b),\qquad
 \mu^{\mathbb R}=\tfrac14(|a|^2-|b|^2).
 \label{eq:e2-double-cover}
\end{equation}
Equal absolute values again lift every complex $\beta$ at zero
real moment map. The cover coordinates are not chiral-ring
generators. The full ring also contains a nilpotent operator;
its scheme-theoretic role is recorded in Appendix~\ref{app:deformations}.

For the root $R=e_1-e_2$ on $S_7$, the compact charge of the
raising current is the difference of the two exceptional-curve
charges:
\begin{equation}
 q_I(e_1)-q_I(e_2)=D_I\iota_*R.
 \label{eq:e2-charge-difference}
\end{equation}
It is not twice this difference. In the standard $SU(2)$ Cartan
normalization, where a root has weight two, the corresponding
coefficient is $\tfrac12D_I\iota_*R$. The effective action on
the reduced ring need not by itself extend to all states;
Appendix~\ref{app:integral-charges} checks the extension in $X_9$
using the diagonal center identification with the other Abelian
and Coulomb factors.

The degree-six del Pezzo cone similarly has $E_3$ flavor algebra
$\mathfrak{su}(3)\oplus\mathfrak{su}(2)$. Its two reduced
deformation components are a Cartan plane and a Cartan line,
corresponding to the $A_2$ and $A_1$ root sublattices of
$K_{S_6}^\perp$. A \emph{branch profile} $\sigma$ chooses one
component at each such point. No $E_3$ point occurs in $X_9$;
the branchwise formulation is useful for stating the general
comparison precisely.

\subsection{The rigid and geometric constraints}

Let $\mathcal B_\sigma$ be the direct sum of the nodal lines,
the reduced $E_2$ lines and the selected $E_3$ Cartan spaces.
Choose compatible root-marked coordinates. The map
\begin{equation}
 Q_{X,\sigma}:\mathcal B_\sigma\longrightarrow
 H^2(\wideX;\QQ)^\vee\otimes\CC
 \label{eq:compact-charge-map}
\end{equation}
has columns the divisor pairings with exceptional nodal curves
and the pushforwards of the selected roots. Since moment maps
add under diagonal gauging, its components are the total complex
moment maps of the rigid local-sector model
\cite{Ceresole:2000}.

\begin{proposition}\label{prop:physical-kernel}
The image of the zero-level rigid triplet-moment-map quotient
in the invariant coordinates $\mathcal B_\sigma$ is
$\ker Q_{X,\sigma}$.
\end{proposition}
\begin{proof}
The complex moment-map equations give one inclusion. Conversely,
lift each coordinate of a vector in the kernel by the balanced
representatives above. For an $E_3$ component use the rank-one
ADHM representative with $|q_i|=|\widetilde q_i|$ and
$\sum_iq_i\widetilde q_i=0$, as detailed in
Appendix~\ref{app:deformations}. Each local real Cartan moment
map vanishes separately, and the complex equations vanish by
the kernel condition. Quotienting does not change these invariant
coordinate values.
\end{proof}

The same kernel has a geometric interpretation. We use the fan
polytope convention: $\mathbb P_\Delta$ is defined by cones over
faces of a reflexive $\Delta\subset\ZZ^4\otimes\mathbb R$;
the polar $\Delta^\circ$ is its anticanonical Newton polytope.
Call $\Delta$ admissible when every edge has lattice length one,
its two-faces are unimodular triangles, unit squares or unit-edge
pentagons/hexagons with one interior lattice point, and every dual
edge of a singular two-face has length one. A $\Delta$-regular
anticanonical hypersurface then has exactly the isolated local
germs just described. These hypotheses make the ambient divisor
calculation detect the full rational kernel
\cite{Batyrev:1994,Wuebben:2026obstruction}.

\begin{theorem}\label{thm:first-order-image}
Let $X$ be such a hypersurface, with $\omega_X\cong\mathcal O_X$,
$h^1(X,\mathcal O_X)=0$, and a projective crepant resolution.
Identifying each branch tangent space with its root-marked
invariant coordinates, the image of the global first-order
deformation space $\mathbb T_X^1$ satisfies
\begin{equation}
 \operatorname{im}\left(\mathbb T_X^1\longrightarrow
                     \bigoplus_pT^1_{X,p}\right)
       \cap\mathcal B_\sigma=\ker Q_{X,\sigma}.
 \label{eq:first-order-image-theorem}
\end{equation}
\end{theorem}
The proof in Appendix~\ref{app:deformations} identifies the
local-to-global obstruction map with the pushforward of link
classes. Combined with Proposition~\ref{prop:physical-kernel},
it identifies two images in $\mathcal B_\sigma$; it does not
identify the entire Higgs branch with a complex vector space.

The exact mixed smoothing theorem strengthens this tangent comparison
\cite[Theorems~4.5 and~4.7]{Wuebben:2026smoothing}. Let
$\mathcal D_\sigma$ be the union of the local discriminants in
$\mathcal B_\sigma$. Then
\begin{equation}
 X\text{ smooths on }\sigma
 \quad\Longleftrightarrow\quad
 \ker Q_{X,\sigma}\setminus\mathcal D_\sigma\ne\varnothing.
 \label{eq:exact-smoothing-criterion}
\end{equation}
At every rank-one summand the coefficient must be nonzero; on an
$E_3$ plane all three coefficients of its root-marked vector
$\mu_1e_1+\mu_2e_2+\mu_3e_3$, with $\sum_i\mu_i=0$, must be
nonzero. Necessity includes arbitrary analytic smoothing arcs.
The selected deformation base is smooth if the relation space projects
nontrivially to every selected $E_3$ summand, and every tangent on it
integrates. Appendix~\ref{app:deformations} states the precise scope.
For $X_9$ there is no $E_3$ hypothesis to check, and its one-dimensional
kernel has a vector nonzero at all three germs. The explicit smoothing
in Section~\ref{subsec:x9} additionally identifies the smooth phase
needed for the charge and period comparisons.

\subsection{The finite-Planck qualification}
\label{subsec:supergravity-qualification}

The charges are exact compactification data, whereas
Proposition~\ref{prop:physical-kernel} uses rigid hyperk\"ahler
cones. A compact supergravity hypermultiplet manifold is instead
quaternionic K\"ahler. Consider a two-derivative description with
Abelian hypermultiplet-isometry gaugings, neutral vector scalars
and no tensor multiplets. Let $P_I^r$ be the quaternionic moment
maps, $k_I^X$ the Killing vectors, and $h^I$ the very-special
coordinates normalized by $\kappa_{IJK}h^Ih^Jh^K=6$.
At constant scalars and zero field strengths, the fermion shifts
give the fully supersymmetric Minkowski conditions
\cite{Ceresole:2000,Mohaupt:2004de}
\begin{equation}
 h^IP_I^r=0,\qquad h_x^IP_I^r=0,\qquad h^Ik_I^X=0.
 \label{eq:compact-susy-shifts}
\end{equation}
At a regular point, $h^I$ and its independent tangent vectors
$h_x^I$ span the full vector space: the cubic gradient annihilates
the tangents but has nonzero pairing with $h^I$. Hence
\begin{equation}
 P_I^r=0\ \hbox{for all }I,r,\qquad h^Ik_I^X=0.
 \label{eq:compact-quaternionic-vacua}
\end{equation}
Gravity does not replace these individual equations by only the
graviphoton combination. On the relevant massless locus the
hyperino condition can vanish for the transition-sector charges;
the quaternionic moment-map equations remain.

The charge embedding alone does not construct the nonlinear map
from $P_I^r$ to the rigid coordinates or a smooth hypermultiplet
action at the interacting $E_2$ point. Explicit conifold
supergravity models already approximate this metric
\cite{Mohaupt:2004de}. We therefore keep three statements separate:
the rigid invariant-coordinate constraint, the compact geometric
smoothing, and the exact vector-sector comparison. The first two
do not supply the missing quaternionic dictionary.
